\section{Monitoring, Telemetry, and Incident Response}

\subsection{Telemetry Stack}

部署的 Agent 每一次工具调用、每一条 API 输出都会被写入观测链路，流向三类系统：

\begin{itemize}
    \item \textbf{Alerting}：基于 anomaly detection 设置阈值，如突发的自动化部署或高频 requests
    \item \textbf{Replay}：完整记录行为用于复现，包含 timeline、embeddings、输入/输出
    \item \textbf{Metrics}：自定义指标（alignment drift、policy deviation、human override rate）
\end{itemize}

\begin{warningbox}{Observability isn't optional}
一旦 Agent 脱离常规模式，必须能够回放完整的调用链、看到 system prompt 以及外部 context。没有这个可视化，调查和治理几乎不可能。
\end{warningbox}

\subsection{Incident Response Loop}

当监控触发警报时，Anthropic 的响应流程包括：

\begin{enumerate}
    \item 快速判定风险等级（green/orange/red）
    \item 自动切换至木偶式模式，所有输出需人工审批
    \item 开启 post-mortem，记录 timeline + root cause
    \item 更新 checklist 与 prompt 指令，防止 recurrence
\end{enumerate}

\subsection{本章小结}

Agent 监控需要 telemetry、alert、replay 三条腿支撑；一旦触发警报，组织必须以 structured loop（等级判定、containment、post-mortem、update）迅速响应，防止能力滑向风险区。

\subsection{Data Pipelines and Instrumentation}

Telemetry data is ingested into a multi-tier pipeline: real-time streaming, batch aggregation, and backing datastore for audit. Each agent interaction emits structured events with the following fields:

\begin{itemize}
    \item `task\_id`, `agent\_variant`, `score`
    \item `tools\_used` list with duration/evidence per tool
    \item `alignment\_drift` delta vs baseline
    \item `human\_override` flag
\end{itemize}

\begin{importantbox}{Why instrumentation matters}
Instrumentation not only powers alerts but also enables retrospective capability audits: you can slice by agent version, time window, tool set, or base capability and replay the exact logs that triggered a red flag.
\end{importantbox}

The ingestion stack keeps both raw and derived layers. Raw events are immutable and tagged with `event\_hash` to ensure idempotency. Derived views pre-aggregate per ASL level and feed dashboards that teams monitor in daily standups.

Dashboards themselves include `alignment drift`, `tool success`, `red-team exposure`, and `human override latency`. When any metric crosses thresholds, the telemetry system automatically tickets the incident response crew.

